\scanpage{145}
\mspage{V.17}
% archive sheet number: 144

Finalement, on trouve un diagramme commutatif \deleted{à carrés cartésiens}
\diagram{\begin{tikzcd}[ampersand replacement=\&,column sep=large,row sep=large]
\text{(4.2)} \& V_s\arrow[r,hook,"i"] \& V \& V_{\unsure{\eta|y}}\arrow[l,hook',"j"'] \\
\& V_{\bar s}\arrow[u,"h_s"]\arrow[r,hook,"\bar i"'] \& \bar V\arrow[u,"h"] \& V_{\bar\eta}\arrow[u,"h_\eta"']\arrow[l,hook',"\bar j"]
\end{tikzcd}},
dans lequel \deleted{les flèches verticales} (à isomorphisme de topos étales près tout au moins, en ce qui concerne $V_{\bar s}$) $V_{\bar s}$ et $V_{\bar\eta}$ sont resp. les images inverses des sous-schémas \inserted{fermé resp. ouvert} $V_s$ et $V_\eta$ de $V$, \deleted{et} et où $h_s$ \deleted{est} \inserted{définit} une équivalence de topos étales.

3.5. Les \deleted{\unsure{rôles des}} cinq schémas \inserted{essentiellement distincts} $V_s,V,V_\eta,\deleted{\unsure{V_{\bar s}}}V_{\bar\eta},\bar V$ jouent, du point de vue topologique, des rôles essentiels distincts :

\deleted{$\unsure{V_s|V_0}$} -- $V_s$ définit la structure topologique locale de la fibre $X_0$ en $x$

-- $V$ définit la structure topologique locale du schéma ambiant $X$ en $x$.

-- $V_\eta=V-V_s$ \deleted{\unsure{sert à relier}} \inserted{fournit, un lien entre,} les structures cohomologiques des schémas précédents $V_s$ et $V$

-- $V_{\bar\eta}$ fournit les \deleted{\unsure{cohomologies}} groupes des \deleted{\unsure{cohomologies}} cycles évanescents de $X$ en $x$. [\inserted{Il est} \deleted{\unsure{et est}} relié à $V_\eta$ $V_{\bar\eta}=V_\eta\times_\eta\bar\eta$, et sa cohomologie à celle de $V_\eta$ par la suite \deleted{exacte} spectrale habituelle].
\margin{par (4.2.1)\[
\left\{\begin{aligned}
(\underline\Phi^i_{\unsure{\nu|v|*}})_{\bar x}&\simeq H^{i-1}(V_{\bar\eta},\Lambda_\nu)\quad\text{si }i\geq2\\
0\longrightarrow(\underline\Phi^0_\nu)_{\bar x}&\longrightarrow\Lambda_\nu\longrightarrow H^0(V_{\unsure{\bar\eta|\eta}},\Lambda_\nu)\longrightarrow(\underline\Phi^1_\nu)_{\bar x}\longrightarrow0
\end{aligned}\right.
\]
suite exacte}

-- $\bar V$ permet de mettre en \unsure{relations|relation} les structures cohomologiques de $V_s\simeq V_{\bar s}$, et celle de $V_{\bar\eta}$.

On voit donc que les rôles de $V_\eta$ et $\bar V$
